% Generated by Listas/make_statements.py from pratica-lecture-02.md, item 5. Do not edit.
\textbf{Source.} Muñoz-Garcia \& Toro-Gonzalez (2019), \emph{Strategy and Game Theory: Practice Exercises with Answers}, Ch. 2, Exercise 7 ``Cournot Game of Quantity Competition'', p.~48 {[}PDF p.~61{]}. Answer in the source: yes, a worked answer follows the exercise. Reproduced verbatim (statement only; the book's worked answer follows it).



\begin{quote}

Consider an industry with two firms competing in quantities, i.e., Cournot competition. For simplicity, assume that firms are symmetric in costs, \(c > 0\), with no fixed costs and that they face a linear inverse demand \(p(Q) = a - bQ\), where \(a > c\), \(b > 0\), and \(Q\) denotes aggregate output. Note that the assumption \(a > c\) implies that the highest willingness to pay for the first unit is larger than the marginal cost that firms must incur in order to produce the first unit, thus indicating that a positive production level is profitable in this industry. If firms simultaneously and independently select their output level, \(q_1\) and \(q_2\), find the Nash Equilibrium (NE) of the Cournot game of quantity competition.

\end{quote}

